% Part: lambda-calculus
% Chapter: syntax
% Section: unique-readability

\documentclass[../../../include/open-logic-section]{subfiles}

\begin{document}

\olfileid{lam}{syn}{unq}
\olsection{અનન્ય વાંચનીયતા}

દરેક પદને રચવાની રીત અનન્ય છે કે કેમ એવો પ્રશ્ન થાય; તેનો
જવાબ હા છે. દરેક લૅમ્બડા પદને રચવાની અને તેનો અર્થ સમજવાની
રીત એક જ છે. આગળની ચર્ચામાં, \emph{રચનાક્રમ} એટલે
\olref[trm]{defn:term}ના રચના-નિયમો એક કે અનેક વાર વાપરીને
પદ રચવાની પ્રક્રિયા.

\begin{lem}\ollabel{lem:term-start}
પદની શરૂઆત ચલથી અથવા કૌંસથી થાય છે.
\end{lem}

\begin{proof}
\olref[trm]{defn:term} મુજબ રચાયું હોય તો જ કંઈક પદ ગણાય.
તે \olref[trm]{defn:term-var}નું પરિણામ હોય તો ચલ જ હોવું જોઈએ.
તે \olref[trm]{defn:term-abs} અથવા
\olref[trm]{defn:term-app}નું પરિણામ હોય તો તેની શરૂઆત કૌંસથી
થાય છે.
\end{proof}

\begin{lem}\ollabel{lem:app-start}
અનુપ્રયોગના પરિણામની શરૂઆત બે ખુલતા કૌંસથી અથવા ખુલતા કૌંસ
પછી ચલથી થાય છે.
\end{lem}

\begin{proof}
$M$ અનુપ્રયોગનું પરિણામ હોય તો તે~$(PQ)$ સ્વરૂપનું છે;
તેથી તેની શરૂઆત કૌંસથી થાય છે. $P$ પદ હોવાથી,
\olref{lem:term-start} મુજબ તેની શરૂઆત કૌંસથી અથવા ચલથી
થાય છે.
\end{proof}

\begin{lem}\ollabel{lem:initial}
પદનો કોઈ યોગ્ય આરંભિક ખંડ પોતે પદ નથી.
\end{lem}

\begin{prob}
પદોની લંબાઈ પર આગમનથી
\olref[lam][syn][unq]{lem:initial} સાબિત કરો.
\end{prob}

\begin{prop}[અનન્ય વાંચનીયતા] \ollabel{prop:unq}
દરેક પદનો રચનાક્રમ અનન્ય છે. બીજા શબ્દોમાં, કોઈ રચનાક્રમથી
પદ~$M$ રચાયું હોય તો એ પદ રચી શકે એવો બીજો રચનાક્રમ નથી.
\end{prop}

\begin{proof}
  પદોના રચનાક્રમ પર આગમનથી આ સાબિત કરીએ.
  \begin{enumerate}
    \item $M$ એ~$x$ સ્વરૂપનું છે, જ્યાં~$x$ કોઈ ચલ છે.
      અમૂર્તીકરણ અને અનુપ્રયોગના પરિણામો હંમેશા કૌંસથી શરૂ
      થાય છે; તેથી~$M$ રચવામાં તે વપરાયા હોઈ શકે નહીં.
      આમ~$M$નો રચનાક્રમ
      \olref[trm]{defn:term}\olref[trm]{defn:term-var}નું
      એક જ પગલું હોવો જોઈએ.
    \item $M$ એ~$(\lambd[x][N])$ સ્વરૂપનું છે, જ્યાં~$x$
      કોઈ ચલ અને~$N$ પદ છે. તે એકલો ચલ નથી; તેથી
      \olref[trm]{defn:term}\olref[trm]{defn:term-var}
      મુજબ રચાયું હોઈ શકે નહીં. \olref{lem:app-start} મુજબ
      તે અનુપ્રયોગનું પરિણામ પણ નથી. આમ~$M$ માત્ર~$N$ના
      અમૂર્તીકરણથી જ રચાઈ શકે છે. આગમનની ધારણાથી~$N$નો
      રચનાક્રમ પોતે અનન્ય છે.
    \item $M$ એ~$(PQ)$ સ્વરૂપનું છે, જ્યાં~$P$ અને~$Q$
      પદો છે. તેની શરૂઆત કૌંસથી થાય છે; તેથી તે
      \olref[trm]{defn:term}\olref[trm]{defn:term-var}
      મુજબ પણ રચાઈ શકે નહીં. \olref{lem:term-start} મુજબ
      $P$ની શરૂઆત~$\lambd$થી થઈ શકતી નથી; તેથી~$(PQ)$
      અમૂર્તીકરણનું પરિણામ પણ હોઈ શકે નહીં. હવે ધારો કે
      $M$ને અનુપ્રયોગથી રચવાની બીજી રીત હોય, એટલે કે તે
      $(P'Q')$ સ્વરૂપનું પણ હોય. તો~$P$ એ~$P'$નો યોગ્ય
      આરંભિક ખંડ હશે (અથવા તેનો ઊલટો); પરંતુ
      \olref{lem:initial} મુજબ આ અશક્ય છે. તેથી~$P$ અને
      $Q$ અનન્ય રીતે નક્કી થાય છે, અને આગમનની ધારણાથી
      $P$ અને~$Q$ બંનેના રચનાક્રમ પણ અનન્ય છે.
  \end{enumerate}
\end{proof}

ઉપરના પ્રસ્તાવને વધુ સહેલાઈથી વાંચી શકાય એમ આ રીતે કહી શકાય:
\begin{prop}
  પદ~$M$ નીચેના સ્વરૂપોમાંથી માત્ર એકનું હોઈ શકે છે:
  \begin{enumerate}
    \item $x$, જ્યાં~$x$ એ~$M$થી અનન્ય રીતે નક્કી થતો ચલ છે.
    \item $(\lambd[x][N])$, જ્યાં~$x$ ચલ અને~$N$ બીજું
      પદ છે; બંને~$M$થી અનન્ય રીતે નક્કી થાય છે.
    \item $(PQ)$, જ્યાં~$P$ અને~$Q$ એ~$M$થી અનન્ય રીતે
      નક્કી થતાં બે પદો છે.
  \end{enumerate}
\end{prop}

\end{document}
